\scanpage{116}
\mspage{103}
% archive sheet number: 115

\unsure{comportent|compatibles} bien \deleted{\unsure{par rapport}} \unsure{à} l'opération
\deleted{\unsure{tant aux|tendent à la} formation des ensembles}
\[
\Phi\mathrel{\unsure{\vdash\joinrel\dashv|\longmapsto}}\Phi^{\sim},
\]
quand $\Phi$ est une partie de $\mathrm{Fl}\deleted{\unsure{I}}M(I)$
resp de $\mathrm{Fl}\deleted{\unsure{I}}M(I_{0})$. \inserted{qui est $\unsure{Q|\mathcal{U}}$-\uline{accessible}, i.e.}
telle qu'il existe une \uline{petite} partie
$\Phi_{0}\subset\Phi$ \unsure{avec} $\Phi\subset\Phi_{0}^{\sim}$,
\inserted{plus précisément \unsure{---|...}}
$u_{!}(\Phi^{\sim})\subset(u_{!}(\Phi))^{\sim}$.

Ceci noté, on a par b
\[
 C(I)=(v_{!}(C(I_{0})))^{\sim}
\]
d'où
\[
 u_{!}(C(I))\subset
 \bigl(\underbrace{u_{!}\bigl(v_{!}(C(I_{0}))\bigr)}_{(uv)_{!}(C(I_{0}))}\bigr)^{\sim},
\]
et la deuxième $=$ s'écrit aussi
\[
 (v'_{!}u_{\deleted{\unsure{0}}0!}(C(I_{0})))^{\sim}
 \overset{?}{\subset}(v'_{!}(I'_{0}))^{\sim}=C(I').
\]
Donc tout revient à montrer que
\[
 u_{0!}(C(I_{0}))\subset C(I'_{0}),\tag*{(6.1.16)}
\]
ce qui résulte aussitôt du \unsure{calcul}
de $u_{0!}$ à coups de sommes directes,
et du fait que $C$ est stable par
sommes directes quelconques. Cela
prouve donc $u_{!}(C(I))\subset C(I')$, et on
prouve de $=$ $u_{!}(TC(I))\subset TC(I')$,
\deleted{\unsure{en se réduisant}} \inserted{qui ramène de $=$ au cas particulier} évident
\[
 u_{0!}(TC(I_{0}))\subset TC(I'_{0}).\tag*{(6.1.17)}
\]

